% Part: sets-functions-relations
% Chapter: sets
% Section: pairs-and-products

\documentclass[../../../include/open-logic-section]{subfiles}

\begin{document}

\olfileid{sfr}{set}{pai}
\olsection{जोड्या, क्रमित घटकसमूह आणि कार्तीय गुणाकार}

\begin{explain}
संचांतील घटकांना क्रम नसतो, हे विस्तारात्मकतेवरून निष्पन्न होते.
म्हणून क्रम दर्शवायचा असेल, तर आपण \emph{क्रमित जोड्या} $\tuple{x, y}$ वापरतो.
अक्रमित जोडी $\{x, y\}$ मध्ये क्रमाला महत्त्व नसते:
$\{x, y\} = \{y, x\}$. क्रमित जोडीमध्ये मात्र क्रम महत्त्वाचा असतो:
$x
\neq y$ असेल, तर $\tuple{x, y} \neq \tuple{y, x}$.

संच सिद्धांतात क्रमित जोड्यांचा विचार कसा करावा? दोन क्रमित जोड्यांचे
पहिले घटक समान असतील आणि दुसरे घटकही समान असतील तेव्हा आणि केवळ तेव्हाच
त्या जोड्या समान असतात, ही कल्पना जपणे आवश्यक आहे. म्हणजेच:
\[
  \tuple{a, b}= \tuple{c, d}\text{ तेव्हा आणि केवळ तेव्हाच, जेव्हा }a = c \text{ आणि }b=d.
\]
वीनर--कुराटोव्स्की यांच्या व्याख्येने संच सिद्धांतात क्रमित जोड्यांची व्याख्या करता येते.
\end{explain}

\begin{defn}[क्रमित जोडी]\ollabel{wienerkuratowski}
	$\tuple{a, b} = \{\{a\}, \{a, b\}\}$.
\end{defn}

\begin{prob}
	\olref[sfr][set][pai]{wienerkuratowski} वापरून, $\tuple{a,
	b}= \tuple{c, d}$ तेव्हा आणि केवळ तेव्हाच, जेव्हा $a = c$ आणि $b=d$
दोन्ही असतात, हे सिद्ध करा.
\end{prob}

\begin{explain}
क्रमित जोडीची व्याख्या निश्चित केल्यावर तिच्या साहाय्याने आणखी संचांची
व्याख्या करता येते. उदाहरणार्थ, कधीकधी दोनपेक्षा अधिक वस्तूंचे क्रमित
अनुक्रम हवे असतात: \emph{त्रिके} $\tuple{x, y, z}$,
\emph{चतुष्के} $\tuple{x, y, z, u}$ इत्यादी.
त्रिक म्हणजे ज्याचा पहिला घटक स्वतःच क्रमित जोडी आहे अशी विशिष्ट क्रमित
जोडी असे मानता येते: $\tuple{x, y, z}$ म्हणजे $\tuple{\tuple{x, y},z}$.
चतुष्कांसाठीही हेच लागू होते: $\tuple{x,y,z,u}$ म्हणजे
$\tuple{\tuple{\tuple{x,y},z},u}$, आणि असेच पुढे.
सर्वसाधारणपणे आपण \emph{क्रमित $n$-घटकसमूह} $\tuple{x_1, \dots, x_n}$ यांचा उल्लेख करतो.

क्रमित जोड्यांचे, किंवा इतर क्रमित $n$-घटकसमूहांचे, काही संच उपयुक्त ठरतील.
\end{explain}

\begin{defn}[कार्तीय गुणाकार]
$A$ आणि $B$ हे संच दिलेले असतील, तर त्यांचा \emph{कार्तीय गुणाकार}
$A \times B$ याची व्याख्या अशी करतात:
\[
  A \times B = \Setabs{\tuple{x, y}}{x \in A \text{ आणि } y \in B}.
\]
\end{defn}

\begin{ex}
$A = \{0, 1\}$ आणि $B = \{1, a, b\}$ असतील, तर त्यांचा गुणाकार असा आहे:
\[
A \times B = \{ \tuple{0, 1}, \tuple{0, a}, \tuple{0, b},
    \tuple{1, 1}, \tuple{1, a}, \tuple{1, b} \}.
\]
\end{ex}

\begin{ex}
$A$ हा संच असेल, तर $A$ चा स्वतःशी गुणाकार $A \times A$ हा $A^2$
असाही लिहितात. तो $x, y \in A$ असणाऱ्या \emph{सर्व} $\tuple{x, y}$ जोड्यांचा
संच आहे. सर्व $\tuple{x, y, z}$ त्रिकांचा संच $A^3$ असतो, आणि असेच पुढे.
पुढील पुनरावर्ती व्याख्या देता येते:
\begin{align*}
  A^1 & = A\\
  A^{k+1} & = A^k \times A
\end{align*}
\end{ex}

\begin{prob}
$\{1, 2, 3\}^3$ मधील सर्व !!{element}s लिहा.
\end{prob}

\begin{prop}\ollabel{cardnmprod}
$A$ मध्ये $n$ !!{element}s आणि $B$ मध्ये $m$ !!{element}s असतील,
तर $A
\times B$ मध्ये $n\cdot m$ घटक असतात.
\end{prop}

\begin{proof}
$A$ मधील प्रत्येक !!{element} $x$ साठी, $\tuple{x, y} \in A \times B$
या आकाराचे $m$ !!{element}s असतात.
$B_x = \Setabs{\tuple{x, y}}{y
  \in B}$ असे मानू. $x_1 \neq x_2$ असेल तेव्हा $\tuple{x_1, y} \neq
\tuple{x_2, y}$ असल्यामुळे $B_{x_1} \cap B_{x_2} = \emptyset$.
पण $A = \{x_1,
\dots, x_n\}$ असेल, तर $A \times B = B_{x_1} \cup \dots \cup B_{x_n}$;
म्हणून त्यात $n\cdot m$ !!{element}s असतात.

हे चित्ररूपाने पाहण्यासाठी $A \times B$ मधील !!{element}s जाळीच्या स्वरूपात मांडा:
\[
\begin{array}{rcccc}
  B_{x_1} = & \{\tuple{x_1, y_1} & \tuple{x_1, y_2} & \dots & \tuple{x_1, y_m}\}\\
  B_{x_2} = & \{\tuple{x_2, y_1} & \tuple{x_2, y_2} & \dots & \tuple{x_2, y_m}\}\\
  \vdots & & \vdots\\
  B_{x_n} = & \{\tuple{x_n, y_1} & \tuple{x_n, y_2} & \dots & \tuple{x_n, y_m}\}
\end{array}
\]
सर्व $x_i$ वेगवेगळे आहेत आणि सर्व $y_j$ वेगवेगळे आहेत. त्यामुळे या जाळीतील
कोणत्याही दोन जोड्या समान नाहीत, आणि त्यांची संख्या $n\cdot m$ आहे.
\end{proof}

\begin{prob}
$k$ वरील गणिती विगमनाने असे दाखवा की, प्रत्येक $k \ge 1$ साठी,
$A$ मध्ये $n$ !!{element}s असतील, तर $A^k$ मध्ये $n^k$ !!{element}s असतात.
\end{prob}

\begin{ex}
$A$ हा संच असेल, तर $A$ वरील \emph{शब्द} म्हणजे $A$ मधील
!!{element}s यांचा कोणताही अनुक्रम.
असा अनुक्रम म्हणजे $A$ मधील !!{element}s यांचा $n$-घटकसमूह असे मानता येते.
उदाहरणार्थ, $A = \{a, b, c\}$ असेल, तर ``$bac$'' या अनुक्रमाचा
$\tuple{b, a, c}$ हे त्रिक म्हणून विचार करता येतो.
शब्दांना, म्हणजे चिन्हांच्या अनुक्रमांना, संगणकशास्त्रात अत्यंत महत्त्व आहे.
संकेतानुसार, $A$ मधील !!{element}s यांना आपण $1$ लांबीचे अनुक्रम मानतो,
आणि $\emptyset$ याला $0$ लांबीचा अनुक्रम मानतो.
मग $A$ वरील \emph{सर्व} शब्दांचा संच असा आहे:
\[
A^* = \{\emptyset\} \cup A \cup A^2 \cup A^3 \cup \dots
\]
\end{ex}

\end{document}
